\section{Discussion}

\subsection{Relation to work on preprocessing sensitivity}
\label{sec:disc-lipinski}
Lipi\'nski \cite{lipinski2026sensitivity} approached the same problem from the
analysis side: on cadmium, lead and nickel, discarding below-LOQ measurements,
substituting half the LOQ and substituting the full LOQ gave substantially
different pollution patterns, and the recommendation is that studies report the
rule they used. All three yield a number; none can say the data support no
verdict at all. That is our third outcome. The comparison also needs a below-LOQ
flag and a \texttt{procedureLOQValue}; \num{25.5}\% of the record carries neither
(Section~\ref{sec:results-waterbase}), so those rows were excluded. CENSO
counts them as \texttt{UnresolvedObservation}. Metals are also the mild case:
\num{43.8}\% of Annex~I metal samples fall below the quantification limit,
against \num{77.3}\% of organic micropollutant
samples\src{scripts/22\_waterbase\_external.py}, and no comparable study covers
the newly regulated class.

\subsection{The substitution is mandated, which is the point}
\label{sec:disc-mandated}
Article~5(1) of Directive 2009/90/EC does not permit half-limit substitution, it
\emph{requires} it: a measurand below the limit of quantification ``shall be set
to half of the value of the limit of quantification
concerned''~\cite{eu2009qaqc}. Nothing in this paper is therefore a criticism of
the analysts or the authorities who do it, and the finding is not that a bad
convention is widespread. It is narrower and harder: the instrument that mandates
the substitution is the same one that sets the performance criterion the
substitution then obscures, and the record cannot show both at once. A reporter
who complies with Article~5(1) produces the \num{112034} exceedances of
Section~\ref{sec:results-wbabox} that no measurement supports, and a reporter
who does not comply produces a mean nobody can reconstruct. The way out is not a
better rule for choosing a number; it is a representation in which the bound
survives alongside whatever number the law requires to be entered, so that
Article~5(1) and Article~3(3b) can both be honoured and audited. That is what
carrying the limit onto the result buys.

\subsection{Failures that no preprocessing rule can repair}
Three findings survive any substitution rule. Where the quantification limit
exceeds the threshold, no handling of the non-detects can demonstrate
compliance: \num{17.5}\% of assessments against a European standard, and every
assessment for some pesticides (Section~\ref{sec:results-waterbase}).
Substitution only manufactures a verdict whose direction depends on the rule
chosen (Section~\ref{sec:results-wbabox}).

Several thresholds carry preconditions the record cannot satisfy: cadmium's
Annex~I standard is stated over five water-hardness classes, and the
bioavailability-corrected standards need covariates the aggregated release
lacks. A numeric threshold column carries neither the covariates nor any way to
say the standard could not be applied; CENSO returns
\texttt{PreconditionUnmet} rather than a default pass.

A threshold can also be transcribed wrongly, and nothing in a conventional
workflow would surface it. Both regulation packages released here carry a
machine-readable \texttt{transcriptionStatus}. We found the defect in our own
table: five annual averages had been parsed a factor of ten too strict, because
the extraction renders a multi-digit integer with a space between its digits,
and four maximum-allowable cells Annex~I marks ``not applicable'' had been
filled with the next number on the line. Every count against a European
standard was inflated by it. The values are now diffed against a hand-read of
the primary text on every build\src{scripts/99\_audit.py}, so a mis-parse
fails the pipeline instead of propagating --- which is the machine-readable
provenance this section argues for, applied to the argument's own inputs.

\subsection{What twelve years of reform fixed, and what it did not}
\label{sec:disc-reform}
The record shows real progress, and it should be said plainly. The share of
station-years declaring neither a censoring flag nor a limit falls from
\num{31.8}\,\% in \num{2012} to \num{0.0}\,\% in \num{2013} and has stayed there
for \num{12} years. That is not a trend, it is a step: European reporters now
record the limit in force, every time. The representational correction this paper
asks for at the point of measurement is work they have already done once.

It did not make the problem smaller. It made it \emph{visible}, which is a
different achievement, and three insufficiencies survive it.

\emph{The analytical capability did not follow the record-keeping.} Once the
limit is reported it can be tested, and the test has the same answer it had
eighteen years ago: \num{46.3}\,\% of assessments failed the \num{30}\,\%
criterion in \num{2006} and \num{47.7}\,\% in \num{2024}, with the share whose
limit exceeds the standard outright at \num{28.1}\,\% and \num{28.2}\,\%. The
limit is now always written down, and it is still too high.

\emph{The reform reached the metadata, not the value.} Censored rows carrying a
positive number stay between \num{97.2}\,\% and \num{100}\,\% across every year
the record can speak about. What was repaired is the field that says a result is
censored, not the field that still carries a substituted number in its place ---
which is why the flag is what makes the substitution reversible, and why
\num{25.5}\,\% of the historical record remains beyond repair.

\emph{The standards are moving faster than the chemistry.} Of the substances
Directive~(EU) 2026/805 adds with an annual-average standard, \num{33}\,\% place
it below \num{e-3}~\si{\micro\gram\per\litre} against \num{20}\,\% on the
legacy list --- decades in which the measured undecidable share is already
between \num{82} and \num{100}\,\%. Better reporting cannot close a gap that the
threshold is widening.

And one insufficiency is of a different kind altogether, because no analytical
improvement touches it. For cadmium, lead and nickel the standard is defined on a
quantity the reporting schema has no field for: a hardness class, or a
bioavailable concentration requiring hardness, dissolved organic carbon and pH.
That is \num{131068} assessments, \num{18.8}\,\% of every station-year a European
standard reaches (Section~\ref{sec:results-precondition}). A better laboratory
does not help here. Only a schema with somewhere to put the covariates does, and
a vocabulary that can say the standard was not applicable rather than met.

That last insufficiency is a metals problem, and it is worth saying so, because
the others are not. The censoring this paper is about is worst where the
regulation is newest: organic micropollutants are below the quantification limit
in \num{77.3}\,\% of samples against \num{43.8}\,\% for the Annex~I metals,
the substances whose limits most often exceed their own standards are
pyrethroids, neonicotinoids and pharmaceuticals rather than metals, and the
substances Directive~(EU) 2026/805 adds are almost all of that kind. Metals
supply the one failure mode that is representational rather than analytical; the
micropollutants supply the scale.

\subsection{Why the vocabularies lacked this}
SOSA/SSN was built for in-situ sensing, where a detection limit genuinely is a
property of the device and \texttt{ssn-system:DetectionLimit} is correctly a
\texttt{SystemProperty}. Trace organic monitoring analyses a transported sample
in a run with its own calibration and limits. The dominant water ontologies
inherited the sensing model, scoring heavily on sensing terminology and near
zero on sampling; CENSO puts the epistemics on the result.

That inheritance is why the shortfall is not repaired by adding a term.
Section~\ref{sec:results-separation} shows the consequence on one real record:
a limit bound to the sensor is the same fact whether the result it produced was
a measurement or a bound, so it cannot mark the difference between them, and a
vocabulary carrying it separates the two readings no better than one carrying
no limit at all. The distinction has to move with the limit, onto the result.

\subsection{Limitations}
\begin{itemize}
  \item \textbf{The standards are applied counterfactually.} Annex~I as
    consolidated in \num{2026} is applied to a record none of which was
    reported under it (Section~\ref{sec:materials}). No reporter is in breach
    of a rule that did not yet bind it and nothing here should be read as such
    a finding; what is tested is whether monitoring as performed and reported
    could decide the standards now in force --- a property of the data, not of
    the date. Section~\ref{sec:results-dual} reports how far a number moves
    when the package is changed, which is the same question asked sideways.
  \item \textbf{Anything recent rests on eight reporting authorities.} Only
    \num{8} of the \num{37} reporters contribute a post-\num{2015} river
    station-year (Section~\ref{sec:results-era}), so no claim here
    characterises \emph{current European practice}. What the coverage does
    support is narrower and is what the paper asserts: where the
    record-keeping defect was repaired, the reporter repaired it, and the
    analytical defect survived the repair.
  \item \textbf{We audit what was reported, not what was measured.} A missing
    limit may be a reporting omission rather than a laboratory one. That
    matters for blame, not for the argument: either way the record cannot
    support the verdict drawn from it.
  \item \textbf{The graph is a materialisation, not the analysis.} The
    three-valued assessment is computed for every station-year a European
    standard reaches; what is expressed as RDF is a \num{40000}-row reservoir
    sample with a fixed seed, because a pure-Python triple store cannot hold
    the full record. Nothing the paper claims rests on the sample,
    whose method-insufficient share matches the population count.
  \item \textbf{Group standards are excluded from per-substance comparison.}
    A sum standard is not a limit on any of its members, so those rows are
    reported as not individually assessable rather than assessed.
  \item \textbf{\texttt{PossibleExceedance} rests on a legal bound, not a
    reported uncertainty.} \texttt{PossibleExceedance} needs an interval around a
    quantified result. \emph{No} European release reports a per-result measurement
    uncertainty; the disaggregated release holds only
    \texttt{procedureLOQValue} and a below-quantification flag. Article~4(1) of
    Directive 2009/90/EC requires a limit of quantification at or below
    \num{30}\,\% of the standard and an expanded measurement uncertainty of
    \num{50}\,\% or below at coverage factor $k=2$, estimated at the level of
    the standard. We can count the first; WISE-6 provides nowhere to record the
    second. We therefore apply the legal maximum: the widest interval a lawful
    method could have. The construction is the ordinary decision rule of
    conformity assessment~\cite{jcgm2012conformity,eurachem2021compliance},
    in which an interval overlapping a limit yields neither acceptance nor
    rejection and ISO/IEC~17025 requires the rule to be stated~\cite{iso17025};
    what is new here is not the rule but that the outcome it produces is a
    class a query can return. It makes \texttt{PossibleExceedance} operative ---
    \num{6639} assessments --- but it is a \emph{bound}, not an estimate: none
    can be decided from what is reported, though a laboratory beating the legal
    minimum would decide some. A \emph{censored} result whose limit exceeds the
    standard has an interval straddling the threshold, so two subclasses of
    \texttt{IndeterminateCompliance} describe it; Article~3(3b) is decisive by
    the precedence of Section~\ref{sec:methods-vocab} and those assessments
    appear under \texttt{MethodInsufficient}. Both are indeterminate either way,
    so the precedence chooses the reason recorded, not the verdict.
  \item \textbf{One matrix.} River water only; lakes, transitional and coastal
    waters have their own standards and are not analysed here.
  \item \textbf{The vocabulary comparison is of published files.} An ontology
    may be used in ways its published serialisation does not show, and three
    files could not be retrieved; those are reported as unknown, not absent.
\end{itemize}

\subsection{Future work}
\label{sec:discussion-future}

\paragraph{Beyond compliance: attribution.}
The representation supports an attribution chain this paper does not pursue: a
substance's use category mapped to an expected pressure type, then compared
against the land use draining to a station. Assignments would derive from declared
substance--use--sector relations rather than from a factor interpreted after
the fact, and the chain returns ``cannot be determined'' where positive matrix
factorisation over zero-substituted data returns a number. It needs a denser
station network than any survey here provides.

\paragraph{Ontology evaluation.}
All \num{20} competency questions run against the populated graph return
non-empty answers, with their SPARQL
recorded\src{scripts/17\_run\_competency\_questions.py}. The OOPS!\ scanner
reports no critical pitfall on the distributed
ontology\src{scripts/16\_export\_for\_scanners.py}; the one it raised was a
real modelling error, now fixed. The FOOPS!\ FAIR assessment of
Section~\ref{sec:methods-repro} leaves two kinds of work outstanding. Three
missing annotations --- a citation, a publisher, an issue date --- are edits to
the file and will land with the archived release, since two of them are not
determinable until it has a DOI. The other three are absence from a public
registry, which is a submission rather than an edit and which we have not made:
a vocabulary should be indexed where people look for one, and until it is, the
reuse this paper argues for costs a reader more than it should.

\paragraph{The taxonomy is declared but not yet interrogated.}
The four subclasses of \texttt{IndeterminateCompliance} are materialised and
counted (Section~\ref{sec:results-wbabox}), and no competency question asks the
graph to group an indeterminate verdict by its reason. That query is the direct
test of whether the subtyping earns its place, and it is the first addition the
question set should take.

\subsection{Implications}
For monitoring programmes the recommendation costs nothing at the point of
measurement: record with every result the limits of the run that produced it ---
the quantification limit the two provisions are written about, and the detection
limit where the method states one. That alone decides whether censoring can be
reconstructed.

For regulators, an environmental quality standard below the achievable limit of
quantification is unenforceable, and the mismatch is discoverable automatically
once both are machine-readable.

For the semantic web community, a standard's implicit measurement modality
propagates into every extension built on it: until SOSA's sampling vocabulary
is connected to analytical provenance, sample-and-analyse domains will keep
reaching for a sensor model that does not fit.
